\let\negmedspace\undefined
\let\negthickspace\undefined
\documentclass[journal]{IEEEtran}
\usepackage[a5paper, margin=10mm, onecolumn]{geometry}
%\usepackage{lmodern} % Ensure lmodern is loaded for pdflatex
\usepackage{tfrupee} % Include tfrupee package

\setlength{\headheight}{1cm} % Set the height of the header box
\setlength{\headsep}{0mm}     % Set the distance between the header box and the top of the text

\usepackage{gvv-book}
\usepackage{gvv}
\usepackage{cite}
\usepackage{amsmath,amssymb,amsfonts,amsthm}
\usepackage{algorithmic}
\usepackage{graphicx}
\usepackage{textcomp}
\usepackage{xcolor}
\usepackage{txfonts}
\usepackage{listings}
\usepackage{enumitem}
\usepackage{mathtools}
\usepackage{gensymb}
\usepackage{comment}
\usepackage[breaklinks=true]{hyperref}
\usepackage{tkz-euclide} 
\usepackage{listings}
% \usepackage{gvv}                                        
\def\inputGnumericTable{}                                 
\usepackage[latin1]{inputenc}                                
\usepackage{color}                                            
\usepackage{array}                                            
\usepackage{longtable}                                       
\usepackage{calc}                                             
\usepackage{multirow}                                         
\usepackage{hhline}                                           
\usepackage{ifthen}
\usepackage{lscape}
\begin{document}

\bibliographystyle{IEEEtran}



\title{5.9.11}
\author{EE25BTECH11057 - Rushil Shanmukha Srinivas
}
% \maketitle
% \newpage
% \bigskip
{\let\newpage\relax\maketitle}

\renewcommand{\thefigure}{\theenumi}
\renewcommand{\thetable}{\theenumi}
\setlength{\intextsep}{10pt} % Space between text and floats

\numberwithin{equation}{enumi}
\numberwithin{figure}{enumi}
\renewcommand{\thetable}{\theenumi}
\textbf{Question:}
Sumit is 3 times as old as his son. Five years later, he shall be two and a half times as old as his son. How old is Sumit at present ?

\textbf{Solution:}
Let Age of son = x

Age of Father = y =3x
\begin{align}
3x-y=0
\end{align}
Five years later, father shall be two and a half times as old as his son, So
\begin{align}
y+5 = \frac{5}{2}(x+5)
\end{align}
\begin{align}
5x-2y = -15
\end{align}
The system of equations formed is :
\begin{align}
3x-y = 0
\end{align}
\begin{align}
5x-2y = -15
\end{align}
Writing it in the matrix form,

\begin{align}
  \myvec{3 & -1\\5 & -2}\myvec{x\\y} &= \myvec{0\\-15}
\end{align}

Forming the augmented matrix to solve the system of equations,

\begin{align}
  \myaugvec{2}{3 & -1 & 0\\5 & -2 & -15}
\end{align}

Using Gaussian Elimination,

\begin{align}
 \myaugvec{2}{3 & -1 & 0\\5 & -2 & -15} 
 \xleftrightarrow{\;R_2 \to 3R_2 -5R_1}
 \myaugvec{2}{3 & -1 & 0\\0 & -1 & -45} 
\end{align}

Using back substitution we get,

\begin{align}
  -y = -45
  \end{align}
  \begin{align}
  y = 45
  \end{align}
\begin{align}
3x-y = 0
\end{align}
\begin{align}
  3x = 45
  \end{align}
  \begin{align}
  x = 15
\end{align}


The solution for the system of equations is :

\begin{align}
  \myvec{x\\y} &= \myvec{15\\45}
\end{align}

Therefore the present age of father = y = 45

\begin{figure}[h!]
\centering
\includegraphics[width=0.9\columnwidth]{figs/solution.png} 
\caption*{Fig: Representation of Lines and Point of Intersection}
\label{Fig12}
\end{figure}

\end{document}